\subsection{Archives audiovisuelles}

\begin{itemize}


\item Spécificités  

"When you start looking at moving images and doing ML on them, your object of study
is no longer an entire film but shots or scenes or still images. And that suddenly brings
in important questions for archives and libraries. When the record itself changes…that
raises huge questions that folks are not grappling with. Not just enriching a particular
object but changing the object". 
\footcite{cordell}

Dificultés de ce type de contenu : 
"moving image data, which is a particularly difficult data type to atomize."
\footcite{cordell}

\item Indexer les archives audiovisuelles 

La spécificité du langage cinématographique comme obstacle à l'indexation 

archives muettes et anciennes 

Contrairement aux documents textuels qui comportent souvent des éléments d'indexation dans leur forme (couverture de livre, ) les archives audiovisuelles n'en contiennent peu ou pas, ce qui les rend opaques et peu réceptives aux pratiques d'indexation. Il s'agit alors de mettre en place une « indexation par le contenu ». \footcite{zotero-item-582} (aussi placé dans IAcinema)

	Premier outil d'indexation des images : Perzeptron de Franz Rozenblatt en 1957 : machine de reconnaissance d'images.

Pour l'indexation des archives audiovisuelles : apprentissage supervisé, non supervisé ou en profondeur : sont appliqués ou combinés. 

"Dans le cas des archives audiovisuelles, on voit en outre très clairement	comment les besoins des utilisateur·rice·s sont pris en compte lors de la conception de	l'indexation automatique en profondeur et comment les données et les traces des	utilisateur·rice·s sont utilisées pour optimiser l'accès aux archives audiovisuelles (mot-clé :	système de recommandation)."
\footcite{zotero-item-584}

"L'indexation, la recherche et le filtrage basés sur le texte de	contenus audiovisuels pertinents se heurtent par nature à des limites : Les contenus	audiovisuels doivent également pouvoir être recherchés à l'aide de méthodes audiovisuelles	(par exemple la navigation par images, Query by Example)."
\footcite{zotero-item-584}

"L'identification de contenus simples d'images ou de sons basée sur des descripteurs techniques ou des attributs dérivés a tendance à être de plus en plus laissée à la machine apprenante (par exemple « photo de groupe », « personne avec un casque sur la tête », termes parlés), tandis que l'identification plus poussée de personnes, de corps, de lieux, d'événements ou de sujets concrets est encore très sujette à erreurs."
\footcite{zotero-item-584}

\item modification de la perception de l'archive 

Indexer automatiquement une archive audiovisuelles pemret de repenser ses modalités d'analyse et d'interprétation. voir article cairn.

\end{itemize}